
\chapter{Conclusion}
\label{ch:conclusion}

This thesis asked what a small language model running on a phone should do with a conversation it
can no longer fit in its context window. The question was answered by measuring, in one pipeline,
what each memory strategy delivers in answer quality and what it costs on the device that has to
run it.

\section{Findings}
\label{sec:ch10_findings}

\paragraph{Retrieval choices mattered as much as reader size.}
Vector retrieval answered more questions correctly than truncated full context for all five readers
on the phone, had the highest accuracy point estimate for every reader on the workstation under the
same truncation, and kept that ordering on LongMemEval-S. Full context failed because of the
truncation: with a long window, the three readers tested matched or exceeded vector retrieval on
LoCoMo, where they read the whole history, and trailed it, without separating, on LongMemEval-S, with prompts about 39 and 48 times
longer. For readers of 0.8 to 4 billion parameters under a 4,096-token history
budget, the tested range of retrieval depth moved accuracy more than the tested doublings of reader
size. For readers of this size, the evidence placed in front of the reader is therefore worth tuning
alongside the choice of model.

\paragraph{On the phone, the better answer was also the cheaper one.}
Vector retrieval had higher accuracy point estimates, lower energy per answer and shorter query times
than truncated full context in all five reader pairs, with 5.8 to 26.9 times more correct answers per
kilojoule where the ratio is finite. Building a graph in the phone pilot cost about 26 times more
energy per dialogue turn than embedding, a cost paid before the first question that prompt-token
accounting does not show, and peak RAM and temperature did not follow the energy ordering.

\paragraph{The graph lost evidence that its own sources still held.}
On the same retrieved turns, restoring the source text behind the selected triples raised Gemma~4~E2B
from 21.5\% to 33.0\%: adding the associated source text recovered much of the observed deficit. Below
that sits an extraction floor: where most
extraction batches failed, as for the smaller Qwen extractors and the phone pilot, retrieving more
facts did not recover the answers.

\paragraph{The proposed selection rule did not work.}
Selective Provenance Expansion failed its first pre-registered criterion, and its second passed
numerically without a distinguishable gain; the repaired selector chose no source turn in any of the
800 pre-specified prompts. A post hoc oracle using judged labels shows headroom at the question
level, so the negative result bounds this per-fact rule rather than selective expansion itself.

\section{Implications for on-device memory}
\label{sec:ch10_implications}

The most general implication is that answer quality and device cost have to be measured in the same
experiment, on the device, because prompt-level proxies miss the cost of building memory. Three design
consequences follow.

\begin{enumerate}
  \item \textbf{Price the whole lifecycle, not the prompt.} Construction, retrieval, prefill and
        decode are separate operations, and the strategies separated most sharply in construction,
        which never appears in a prompt-token count.

  \item \textbf{Default to bounded retrieval of source text, and keep device limits visible.} On
        this phone it had higher accuracy point estimates, shorter query times and lower energy per
        correct answer than truncated full context. On the workstation it trailed long-window full
        context by up to 8 points on LoCoMo and stayed ahead of it on LongMemEval-S, with 2 to 3\% of the
        prompt tokens. A compact derived representation must earn its construction bill, and peak memory and
        temperature still decide whether a configuration can run at all.

  \item \textbf{Instrument for provenance and repeatability.} Digests, explicit measurement spans
        and checkpoints turn a battery-consuming run into evidence. A budgeted selection rule also
        needs a defined tie-break: in the first graph retrieval rule, 164 of 200 questions had a score
        tie at the fact budget.
\end{enumerate}

\section{Future work}
\label{sec:ch10_future_work}

\paragraph{Close the measured gaps first.}
The most direct next steps address the limits stated in Chapter~\ref{ch:discussion}: a second
annotator and a second judge model for the human audit, adding matched full-context phone sittings so that every
efficiency ratio rests on comparable replication, and running a complete repaired graph-QA cell on
the phone with separate meters for construction and query service. A few whole-history LoCoMo
prompts should also be timed and metered on the phone, so that the cost of long-window full context
rests on a measurement rather than on the arithmetic in Section~\ref{sec:ch4_long_fc}. Because building the memory took
about 60\% of the energy in a 20-question E2B vector sitting, the same questions should also be
repeated at 20, 50 and 100 queries per history, to measure how construction amortizes with reuse.

\paragraph{Cheaper construction.}
Graph memory cost 52.3~J per turn on the phone because a general-purpose small model wrote it. A
specialized extractor, such as the distilled entity recognizer of SmartRAG~\citep{smartrag}, or a
construction pipeline distilled into one compact model, as in MemoryCPT~\citep{memorycpt}, would
attack that term directly. Such a component should be judged by the ledger used here: joules for
construction and joules per correct answer, not tokens alone.

\paragraph{A better representation to store.}
The measured deficit belongs to one representation: short triples produced by a small model, with
their source spans reachable but not present. Representations that keep text with their structure,
or store a retrieval gist beside a pointer to the original turn, should preserve the qualifiers and
temporal detail that source restoration returned, and can be tested against the same matched
comparison.

\paragraph{Decide expansion once per question.}
The SPE analysis points to a utility that is comparable across representations and to a decision
taken once per question instead of once per fact. Reserving part of the budget for source text, or
deciding at query level as LeanMem does with a planner~\citep{leanmem}, are the variants the oracle
bound justifies testing.

\paragraph{Beyond question answering.}
The measured task is memory-dependent question answering. A phone agent also acts: it calls tools,
changes device state and observes the result. Extending this evaluation to tool use means scoring
whether the memory supplied the evidence behind a correct action, at a measured cost.
